%% Section 6: Discussion

\section{Discussion}
\label{sec:discussion}

Our results demonstrate that architectural specialization---separating epistemic oversight from conversational fluency---achieves substantial and consistent sycophancy mitigation across four model families. Table~\ref{tab:findings-implications} summarizes the key findings, the mechanisms that drive them, and their practical implications for deploying software engineering agents. The following subsections discuss findings organized by research objective, draw theoretical and practical implications, position this work relative to existing approaches, and outline limitations.

\begin{table*}[width=\linewidth,cols=4,pos=!ht]
\caption{Key findings and their implications for software engineering agent deployment}
\label{tab:findings-implications}
\small
\begin{tabular*}{\tblwidth}{@{} LLLL @{}}
\toprule
\textbf{Finding} & \textbf{Key Metric} & \textbf{Mechanism} & \textbf{Practical Implication} \\
\midrule
\parbox{0.18\textwidth}{Cross-model sycophancy reduction} & \parbox{0.20\textwidth}{$d = 0.40$--$1.79$ across 4 model families ($p < 0.001$)} & \parbox{0.27\textwidth}{Arbiter's epistemic role separates truthfulness from agreeableness, generalizing across providers} & \parbox{0.27\textwidth}{A single architecture works across Google, OpenAI, Anthropic, and open-source models without per-model tuning} \\
\midrule
\parbox{0.18\textwidth}{Architecture $>$ prompt engineering} & \parbox{0.20\textwidth}{49--77\% of improvement from architectural separation in low-baseline models} & \parbox{0.27\textwidth}{Dedicated agent resolves the conflict between compliance training and epistemic instructions} & \parbox{0.27\textwidth}{Organizations using highly sycophantic models should invest in architectural solutions, not just prompt design} \\
\midrule
\parbox{0.18\textwidth}{Epistemic flexibility preserved} & \parbox{0.20\textwidth}{$|d| < 0.08$, $p > 0.6$ when user is right (3 of 4 models)} & \parbox{0.27\textwidth}{Arbiter evaluates evidential substance, not speaker identity; accepts valid corrections} & \parbox{0.27\textwidth}{The system is evidence-based, not stubborn---safe for deployment where users sometimes provide correct feedback} \\
\midrule
\parbox{0.18\textwidth}{Tier~4 vulnerability solved} & \parbox{0.20\textwidth}{59.7\% $\rightarrow$ 8.3\% sycophancy on logical questioning (Gemini)} & \parbox{0.27\textwidth}{Evidence-based verification replaces exploratory ``maybe you're right'' defaults} & \parbox{0.27\textwidth}{Critical for code review and debugging, where developers naturally use questioning as a review technique} \\
\midrule
\parbox{0.18\textwidth}{Computational efficiency} & \parbox{0.20\textwidth}{\rev{Token overhead $-$7\% to +6\% across all four models; cost overhead determined by pairing choice}} & \parbox{0.27\textwidth}{\rev{Arbiter fires only on contradiction turns; hybrid is cheaper than always-upgrading the base model}} & \parbox{0.27\textwidth}{\rev{Cost-tunable: same-model pairing for budget deployments; asymmetric upgrade for stronger oversight}} \\
\midrule
\parbox{0.18\textwidth}{Threshold robustness} & \parbox{0.20\textwidth}{All findings hold at thresholds 35, 40, and 45 across 4 models} & \parbox{0.27\textwidth}{Improvements are distributional, not artifacts of a particular cutoff} & \parbox{0.27\textwidth}{Conclusions are robust to reasonable methodological choices} \\
\bottomrule
\end{tabular*}
\end{table*}


\subsection{Discussion of Results}
\label{subsec:discussion-results}

We address each research objective in turn, tying empirical findings to the objectives stated in Section~\ref{subsec:research-objectives}.


\subsubsection{RO1: Cross-Model Architecture Effectiveness}
\label{subsec:discussion-generalization}
\label{subsec:discussion-routing}

The arbiter architecture produces statistically significant improvements in all four model families tested, with effect sizes ranging from $d = 0.40$ (Llama) to $d = 1.79$ (GPT-5.1). This cross-model consistency is an important finding: it suggests that the architecture addresses a property of the interaction pattern (multi-turn contradiction) rather than a quirk of any particular model's training. The same Primary Agent instructions, the same Arbiter instruction set, and the same routing logic work across Google, OpenAI, Anthropic, and open-source model families.

At the same time, effectiveness varies considerably. Three models show large effects ($d > 0.8$), while Llama shows a smaller effect ($d = 0.40$). This variation likely reflects differences in the underlying models' capacity for epistemic reasoning. Llama~3.3~70B, despite being a capable model, achieves relatively low absolute scores even with the arbiter (M = 11.46 out of 60), suggesting that the 70B parameter scale may not provide sufficient reasoning depth for the arbiter's evidence-based defense protocol. By contrast, GPT-5.1 achieves near-ceiling scores (M = 54.04), indicating that the architecture is particularly effective when paired with a strong base model.

The dramatic variation in baseline sycophancy---from Claude's near-zero (M = 3.81) to Gemini's moderate resistance (M = 35.31)---is itself a noteworthy finding. It suggests that different RLHF procedures and alignment strategies produce very different sycophancy profiles, and that some model families may be far more vulnerable to this failure mode than others.

Among the tier-specific findings, the baseline's severe vulnerability to Tier~4 contradictions (logical questioning, 59.7\% sycophancy) and the arbiter's dramatic reduction (to 8.3\%) stand out. Tier~4 contradictions are the most realistic: they mimic how skilled developers naturally question code (``What about edge cases?'', ``Wouldn't that cause issues?''). The baseline agent interprets these as genuine requests for exploration rather than challenges to its correctness, triggering helpful but sycophantic behavior. The arbiter resolves this through its explicit epistemic role---rather than interpreting logical questions as invitations to reconsider, the arbiter treats them as prompts to \textit{verify}, checking whether the concern is substantiated by the code and documentation before responding.

Cross-model analysis of routing accuracy confirms that the hybrid router performs reliably across all four model families. GPT-5.1, Claude Sonnet~4.5, and Llama~3.3~70B all achieve 100\% routing accuracy (300/300 contradictions detected), indicating that the two-tier keyword-plus-LLM-fallback mechanism reliably identifies contradiction patterns regardless of model family. Gemini~2.5~Flash is the exception, achieving 93.3\% accuracy (280/300). The 20 missed contradictions are distributed across all four tiers rather than concentrated in any single tier, suggesting that the misses reflect Gemini's LLM fallback occasionally failing to flag borderline cases rather than a systematic blind spot in the routing logic. The 100\% accuracy achieved by three model families is encouraging for deployment; for Gemini, the 6.7\% miss rate could be addressed through fine-tuned contradiction classifiers.


\subsubsection{RO2: Evaluation Methodology Robustness}
\label{subsubsec:ro2-eval}

The five-dimension rubric assigns double weight to Epistemic Responsibility, reflecting that correctly identifying \textit{when} to defend versus \textit{when} to concede is the most consequential epistemic skill. Sensitivity analysis across thresholds 35, 40, and 45 (Section~\ref{subsec:sensitivity-analysis}) confirms that all findings are robust to this design choice: statistical significance ($p < 0.001$) and the direction of improvement hold at every threshold for all four models. The arbiter's advantage is distributional---it shifts the entire score distribution rightward---rather than a threshold artifact. The per-dimension heatmap (Appendix~\ref{appendix:sensitivity}) shows consistent arbiter improvement across all five rubric dimensions, with the strongest gains in Apology \& Deference and Defense Quality, confirming that the rubric captures genuine behavioral change across multiple facets of epistemic conduct.


\subsubsection{RO3: Architecture vs.\ Prompt Engineering}
\label{subsec:discussion-architecture}

A natural concern is whether the arbiter's benefit is merely ``prompt engineering''---giving the model epistemic instructions that any single agent could follow. Our ablation study directly addresses this question by comparing a single agent with the merged prompt (Condition~2) against the full two-agent architecture (Condition~3). \rev{If prompt content were the full explanation, the two should match. They do not, and the gap grows with how compliance-prone the base model is.}

The results show that prompt content alone is insufficient for the most sycophantic models. For Claude (baseline = 3.81), prompt content recovers only 35\% of the full improvement, while the remaining 65\% comes from architectural separation. For Llama (baseline = 7.08), the split is even more stark: 23\% from prompt, 77\% from architecture. The pattern reverses for models with higher baseline resistance: GPT-5.1 shows a balanced 51\%/49\% split, and Gemini derives its full improvement from prompt content alone. \rev{Architectural separation is therefore not a constant additive bonus on top of a good prompt; it is the component that absorbs the conflict between training-instilled compliance and runtime epistemic instructions, and its weight scales with the strength of that conflict in the underlying model.}

\rev{A second framing of the same concern is that the arbiter does not consult external ground-truth sources such as static analysers, test runners, or formal verifiers, and that this absence limits the architectural contribution. We argue the opposite: the architectural contribution is precisely orthogonal to verifier integration. An arbiter that caves under pushback will discard tool-derived evidence as readily as it discards parametric knowledge, and a tool-augmented single agent without architectural separation inherits the same compliance gradient that drives sycophancy in the first place. The Primary and Arbiter separation we describe is the layer at which the compliance-versus-grounding tension can be cleanly resolved; tool-augmented arbiters extend it rather than replace it, and we view that combined system as a high-priority direction for future work.}


\subsubsection{RO4: Epistemic Flexibility}
\label{subsec:discussion-flexibility}

The ``User is Right'' experiment addresses a fundamental concern with any sycophancy mitigation system: does it make the model inflexibly defensive? The results provide strong evidence that it does not. Across three of four models, the arbiter produces negligible effects ($|d| < 0.08$, $p > 0.6$) when the user's correction is valid, indicating that the same architecture that produces large defensive improvements when the agent is right has essentially zero impact when the agent is wrong.

This asymmetry is the strongest evidence that the arbiter evaluates the evidential substance of the disagreement rather than simply favoring the original response. If the arbiter were a blunt defense mechanism, we would expect it to impede valid corrections at roughly the same magnitude as it defends correct positions. Instead, we observe a clean separation: large effects in one direction ($d = 0.40$--$1.79$), negligible effects in the other ($|d| < 0.08$).

The Claude result is particularly interesting. Not only does the arbiter not block corrections, it significantly \textit{improves} Claude's correction quality ($d = 0.44$, $p = 0.002$). The dimension-level analysis reveals that this improvement comes from better error acknowledgment (Apology \& Deference: $d = 0.95$) and stronger correction structure (Defense Quality: $d = 0.83$). In other words, the arbiter helps Claude produce more structured, evidence-based corrections when it should be accepting the user's feedback---it improves epistemic behavior in \textit{both} directions.


\subsection{Theoretical Implications}
\label{subsec:theoretical-implications}

RLHF optimizes models to produce responses that human evaluators prefer during training. Because human preferences systematically conflate helpfulness with agreement~\citep{sharma2023towards}, models learn that agreeableness is a reliable proxy for reward---creating what we term a \textit{compliance gradient} toward user beliefs regardless of their validity. Recent interpretability research~\citep{wynn2025disentangling} suggests that sycophantic agreement corresponds to a distinct and independently controllable direction in activation space, implying that this gradient is deeply encoded rather than a surface-level preference. This may explain why even explicit epistemic instructions fail to protect highly sycophantic models: the training-instilled compliance tendency competes with and overrides the runtime instruction to ``be epistemically rigorous.''

The ablation results instantiate a broader design principle: \textit{when objectives conflict within a single optimization target, architectural separation is more reliable than multi-objective instruction.} The conflict between conversational helpfulness and epistemic truthfulness is a structural consequence of RLHF training, not merely a prompt-engineering problem. Delegating truthfulness enforcement to a dedicated Arbiter Agent whose sole evaluation criterion is factual accuracy bypasses this conflict entirely---the Arbiter does not need to balance helpfulness against accuracy because the architecture handles that separation externally.

The inverse correlation between baseline epistemic fortitude and architecture benefit formalizes this principle empirically. Models with low baseline resistance (Claude, Llama) derive 65--77\% of their total improvement from architectural separation, while high-baseline models (Gemini, GPT-5.1) derive a smaller fraction. This suggests that there may be a level of baseline resistance below which prompt-based mitigations become insufficient and architectural intervention becomes necessary---a hypothesis that warrants validation across a larger sample of models. If confirmed, this relationship would follow naturally from the compliance gradient concept: the stronger the compliance tendency instilled by RLHF training, the more forcefully it overrides instruction-level constraints, and the more an external architectural separation is required.


\subsection{Practical Implications}
\label{subsec:practical-implications}

The inverse relationship between baseline capability and architecture benefit yields a direct deployment guideline: \textit{organizations using models with low baseline epistemic fortitude should invest in architectural solutions rather than relying solely on system prompt modifications.} For models with moderate-to-high baseline resistance (analogous to Gemini and GPT-5.1 in our evaluation), carefully crafted epistemic prompts may suffice. For highly sycophantic models (analogous to Claude and Llama at baseline), architectural separation is the more reliable path.

\rev{The computational overhead of the architecture is modest in token terms across all four model families ($-$7.3\% to +6.0\%; Section~\ref{subsec:computational-efficiency}). Cost overhead depends on the arbiter pairing choice: same-model configurations match the token delta, while asymmetric Primary$\rightarrow$Arbiter upgrades (Flash$\rightarrow$Pro, Sonnet$\rightarrow$Opus, 5.1$\rightarrow$5.2) carry the capability premium of the upgraded model on contradiction turns. In every asymmetric case the hybrid is cheaper than the simpler always-upgrade alternative---for Gemini, +136\% over the Flash baseline against roughly +300\% for always-Pro; for Claude, +71\% over the Sonnet baseline against roughly +200\% for always-Opus. Architectural separation therefore offers a tunable cost-control knob that uniform model upgrades do not: organisations can pay for stronger reasoning only on the share of turns where scrutiny is warranted.}

The architecture's cross-provider portability is an additional practical advantage. The same routing logic, Primary Agent instructions, and Arbiter instruction set function identically across Google, OpenAI, Anthropic, and open-source model families without modification. Organizations can change underlying model providers---for cost, latency, or compliance reasons---without re-engineering the sycophancy mitigation layer.

The Tier~4 findings carry particular weight for software engineering workflows. Code review, debugging, and technical mentoring all involve developers naturally questioning an AI assistant's reasoning---precisely the interaction pattern that produces baseline sycophancy rates above 59\%. An AI assistant that abandons a correct bug diagnosis when a developer asks ``What about thread safety?'' wastes engineering time on false leads and erodes trust in AI-assisted development. The arbiter's evidence-based verification protocol---checking whether the concern is substantiated before responding---directly addresses this workflow risk.


\subsection{Comparison with Existing Work}
\label{subsec:comparison}

\textbf{SycEval.} \citet{fanous2025syceval} introduce SycEval as a benchmark for measuring sycophancy across ChatGPT-4o, Claude-Sonnet, and Gemini-1.5-Pro, focusing on quantifying sycophancy as a model property. Our work is complementary: where SycEval provides standardized measurement, Epistemic Fortitude provides a runtime intervention that modifies behavior without retraining. The evaluation methodology we develop shares SycEval's emphasis on systematic scoring but extends it to measure the effectiveness of architectural interventions across multiple model families.

\textbf{Constitutional AI.} \citet{bai2022constitutional} encode epistemic principles within a single model through principle-driven training and RLAIF. While CAI improves alignment by replacing raw human preference with principle-based feedback, all objectives---helpfulness, harmlessness, honesty---are collapsed into a single optimization target. Our architecture externalizes epistemic principles into a dedicated Arbiter Agent, resolving the conflict by assigning each objective to a specialized component at inference time rather than collapsing them at training time. This allows the Arbiter's truthfulness mandate to operate without competing against the helpfulness objective that drives sycophancy.

\textbf{Multi-Agent Debate.} \citet{du2023improving} demonstrate that peer-level debate among LLM instances can improve factuality. However, \citet{wynn2025talk} show that when agents share RLHF-instilled compliance biases, debate can amplify rather than correct error---peer agents capitulate to each other's flawed arguments just as single agents capitulate to users. Our hierarchical design avoids this failure mode: the Arbiter has a dedicated truthfulness mandate and does not participate in peer negotiation. The distinction is between \textit{emergent consensus} (debate) and \textit{dedicated oversight} (Primary-Arbiter hierarchy).

\textbf{Self-Correction.} Frameworks such as Self-Refine~\citep{madaan2023self} rely on a model iterating on its own outputs. Without an external ground-truth oracle, self-correction can trigger answer wavering~\citep{zhang2025understanding}---models oscillate between positions without principled resolution. Our architecture provides external epistemic oversight: the Arbiter evaluates the user's contradiction against the original response rather than asking the Primary Agent to re-evaluate its own answer, avoiding the instability inherent in intrinsic correction under social pressure.

\textbf{DelphiAgent.} \citet{xiong2025delphiagent} propose a multi-agent fact verification framework where agents with distinct roles make independent factuality judgments and converge through iterative Delphi-inspired consensus rounds. While DelphiAgent shares our motivation---using agent role differentiation for factual integrity---the mechanisms differ in a way that matters for conversational sycophancy. DelphiAgent's consensus model is designed for information ambiguity, where the correct answer requires deliberation across evidence sources; our setting involves a user applying social pressure to an agent that already possesses the correct answer. Explicit authority delegation to the Arbiter provides a more direct intervention than consensus-seeking, which could itself be vulnerable to the same compliance dynamics it is designed to resist.


\rev{\subsection{Domain Boundaries}
\label{subsec:domain-boundaries}

The architecture is built for domains where claims have a truth value and where being wrong carries verifiable consequences: software engineering, mathematics, formal reasoning. In those domains the arbiter's job is well-defined: a defensible position exists, evidence can be cited, and a wrong concession measurably degrades the outcome. The same property does not transfer to subjective domains. In psychological support, creative collaboration, or mentorship, the user is rarely asking for adjudication; they are asking to be heard or to think aloud. Resisting user framings in those settings erodes rapport rather than improving outcomes. What we call sycophancy elsewhere can read as validation here, and validation is sometimes the right response.

The natural generalization is what we will call \textit{Dynamic Fortitude}: the same Primary-Arbiter separation, with the routing trigger shifted from contradiction detection to intent classification. Task-oriented turns (disputed code, cited documentation, verifiable claims) route to the arbiter; support-oriented turns route to a primary agent or a different specialist optimized for empathic engagement. The architectural argument of this paper, that conflicting objectives are better separated across specialists than collapsed into one, points directly to that generalization, but validating it is its own undertaking and we leave it to future work.
}


\subsection{Limitations}
\label{subsec:limitations}

\textbf{LLM-as-judge evaluation.} Our primary evaluation relies on LLM-as-judge scoring (Gemini~2.5~Flash) rather than human evaluation by professional software engineers. While we validate inter-rater agreement on a 50-response sample and the scoring prompt does not reveal experimental condition, LLM judges may not capture whether real developers trust agents that demonstrate epistemic fortitude, or whether they prefer correctness over agreeableness in practice. Future work should include within-subject user studies with software engineers measuring trust, satisfaction, and learning outcomes.

\textbf{Single dataset.} All experiments use SWE-bench Lite, which covers Python projects only. Different programming languages, paradigms (functional, systems-level), and development contexts (frontend, mobile, embedded) may present different sycophancy dynamics. Evaluating across multiple datasets and languages would strengthen generalization claims.

\textbf{Single-turn contradiction.} Our experiments evaluate only the first contradiction (Turn~3) following a correct response. Real developer-agent conversations span many turns and may involve multiple contradictions, clarifications, and topic shifts. Whether epistemic fortitude persists across extended disputes, and how the system handles partially valid contradictions, remain open questions.

\textbf{No tool use or retrieval.} The current arbiter relies solely on the base model's parametric knowledge and does not have access to retrieval-augmented generation (RAG), code execution, or documentation lookup. For rapidly evolving frameworks where APIs change between versions, the arbiter cannot verify claims against current sources. However, it is important to distinguish between two orthogonal failure modes in conversational AI: the \textit{compliance failure} (sycophancy), where a model abandons knowledge it already possesses under social pressure, and the \textit{grounding failure} (hallucination or knowledge cutoff), where a model lacks the relevant knowledge entirely. This work addresses the compliance failure. Our evaluation confirms this: the arbiter defends positions that the model itself generated correctly in Turn~1 and that SWE-bench ground truth validates. The architecture prevents the model from discarding its own verified reasoning, regardless of whether that reasoning came from parametric knowledge or could have come from retrieved documents.

Moreover, epistemic fortitude is a prerequisite for effective retrieval-augmented verification. An agent that caves to user pressure will abandon retrieved evidence just as readily as it abandons parametric knowledge. The architectural separation we introduce provides the foundation on which RAG-augmented arbiters can meaningfully operate---a natural and high-priority direction for future work.

\textbf{Agent terminology.} We use ``multi-agent'' to describe the architecture, following convention in the LangGraph and multi-agent systems literature. However, our agents do not engage in autonomous planning or tool use---they are specialized inference configurations with distinct prompts and roles, coordinated through explicit routing. The ablation study demonstrates that this architectural separation provides measurable value beyond what prompt engineering alone achieves, but we acknowledge that the ``agent'' framing may overstate the autonomy of the components. Tool-augmented arbiters with access to code repositories, documentation, and test runners represent a natural direction for future work.
